\begin{abstract}
Accurate weather forecasts are critical for scheduling and ensuring the reliability of solar generation, yet the underlying structure of forecast errors offers exploitable patterns for grid operators. Using High-Resolution Rapid Refresh (HRRR) forecasts paired with observations across the 896-site PJM solar fleet over 2022–2024, we characterize forecast error behavior across temporal (lead-time, season, year, error growth, autocorrelation) and spatial (cross-site and cross-zone coherence) components. By studying statistics on forecast error, we discovered that forecast error grows as lead time increases. Furthermore this error growth follows a sinusoidal 24-hr pattern that is statistically significant for GHI, relative humidity, U10 wind, and wind direction. We found that errors are highly concentrated and have low varability at early lead times. However, with each increasing lead hour, the varability in errors increases until hitting a plateau around lead hour 12. Autocorrelation (ACF) and partial autocorrelation (PACF) analyses show that errors are highly localized, dominated by lag-1 effects, and shows no statistically significant variations across zones or years. Spatially, error coherence is strongly variable-dependent: synoptic-scale fields (surface pressure, temperature, relative humidity) remain highly correlated across hundreds of kilometers, whereas cloud-driven global horizontal irradiance (GHI) and wind direction decorrelate rapidly within tens of kilometers ($L_{\text{GHI}} \approx \SI{134}{km}$ at solar noon). While spatial coherence jumps dramatically during year 2024, day-block bootstrap resampling demonstrates that yearly variations in decorrelation length fall within sampling uncertainty. These results provide a rigorous, quantitative framework for localized harmonic bias correction and portfolio-level risk diversification strategies.
\end{abstract}